\documentclass[11pt]{article}

\usepackage{amsmath,amssymb,geometry}
\geometry{margin=1in}

\title{Topological Integrity Gravity (TIG):\\
A Spectral-Derived Effective Extension of General Relativity}

\author{Kai}
\date{}

\begin{document}

\maketitle

\begin{abstract}
We propose Topological Integrity Gravity (TIG) as an effective extension of General Relativity motivated by spectral geometry. Starting from a spectral action based on the Dirac operator, we derive a low-energy truncation containing curvature-squared corrections. An additional logarithmic functional introduces nonlinear feedback at high curvature. The resulting field equations admit a scalar degree of freedom and lead to regularized high-curvature behavior. We analyze black hole solutions, cosmological dynamics, and stability conditions, and outline testable predictions.
\end{abstract}

\section{Introduction}

General Relativity (GR) provides an accurate description of gravitational phenomena but admits singular solutions and lacks intrinsic ultraviolet (UV) regularization. Extensions of GR often introduce higher-order curvature invariants or quantum corrections.

In this work, we investigate an approach motivated by spectral geometry, where spacetime structure is encoded in the spectrum of a Dirac operator. The framework leads naturally to curvature-dependent corrections and suggests a stabilizing feedback mechanism at high curvature.

\section{Spectral Framework}

We consider a spectral triple $(\mathcal{A}, \mathcal{H}, D)$, where $D$ is the Dirac operator. The spectral action is defined as:

\begin{equation}
S_{\mathrm{spec}} = \mathrm{Tr}\, f\left(\frac{D}{\Lambda}\right)
\end{equation}

where $f$ is a cutoff function and $\Lambda$ is a characteristic energy scale.

\subsection{Dirac Operator and Geometry}

On a smooth spin manifold, the Dirac operator is:

\begin{equation}
D = i \gamma^\mu \nabla_\mu
\end{equation}

Its square satisfies the Lichnerowicz formula:

\begin{equation}
D^2 = -\nabla^\mu \nabla_\mu + \frac{1}{4} R
\end{equation}

This establishes a direct relation between the spectrum of $D$ and the scalar curvature $R$.

\subsection{Heat-Kernel Expansion}

For large $\Lambda$, the spectral action admits an asymptotic expansion:

\begin{equation}
\mathrm{Tr}\, f(D/\Lambda)
\sim
\sum_n f_n \Lambda^{4-n} a_n(D^2)
\end{equation}

The Seeley–DeWitt coefficients yield:

\begin{align}
a_0 &\sim \int d^4x \sqrt{-g} \\
a_2 &\sim \int d^4x \sqrt{-g} \, R \\
a_4 &\sim \int d^4x \sqrt{-g} \left(R^2 + R_{\mu\nu}R^{\mu\nu} + \cdots \right)
\end{align}

\section{Effective Action}

We consider a low-energy truncation retaining the leading scalar curvature terms:

\begin{equation}
S_{\mathrm{eff}} =
\frac{1}{16\pi G}
\int d^4x \sqrt{-g}
\left(
R + \frac{\alpha}{\Lambda^2} R^2
\right)
+ S_m
\end{equation}

This corresponds to an $f(R)$-type theory.

\section{Integrity Functional}

We extend the spectral action by introducing:

\begin{equation}
\mathcal{I}[D] =
\mathrm{Tr}\,\log\left(1 + \frac{D^2}{\Lambda^2}\right)
\end{equation}

For small $D^2/\Lambda^2$:

\begin{equation}
\log\left(1 + \frac{D^2}{\Lambda^2}\right)
=
\frac{D^2}{\Lambda^2}
-
\frac{D^4}{2\Lambda^4}
+
\cdots
\end{equation}

This generates additional curvature-dependent corrections.

\section{Field Equations}

Variation yields:

\begin{equation}
G_{\mu\nu} + \mathcal{I}_{\mu\nu} = 8\pi G T_{\mu\nu}
\end{equation}

where, in the $R^2$ truncation:

\begin{align}
\mathcal{I}_{\mu\nu}
&=
\frac{2\alpha}{\Lambda^2} R R_{\mu\nu}
-
\frac{\alpha}{2\Lambda^2} R^2 g_{\mu\nu} \\
&+
\frac{2\alpha}{\Lambda^2}
\left(
g_{\mu\nu} \Box R
-
\nabla_\mu \nabla_\nu R
\right)
\end{align}

\section{Trace Equation}

Taking the trace:

\begin{equation}
-R + \frac{6\alpha}{\Lambda^2} \Box R = 8\pi G T
\end{equation}

which can be rewritten as:

\begin{equation}
\Box R - m^2 R = \frac{8\pi G \Lambda^2}{6\alpha} T
\end{equation}

with:

\begin{equation}
m^2 = \frac{\Lambda^2}{6\alpha}
\end{equation}

This corresponds to a propagating scalar degree of freedom.

\section{Black Hole Solutions}

A regularized metric of the form:

\begin{equation}
A(r) = 1 - \frac{2GM r^2}{r^3 + r_0^3}
\end{equation}

illustrates the suppression of curvature divergences in the core region.

\section{Cosmology}

For a flat FRW metric:

\begin{equation}
R = 6(2H^2 + \dot{H})
\end{equation}

The modified Friedmann equation becomes:

\begin{equation}
3H^2 = 8\pi G \rho + \frac{\alpha}{\Lambda^2}\left(\frac{1}{2}R^2 - 6H \dot{R}\right)
\end{equation}

\section{Stability}

The theory is stable under:

\begin{equation}
\alpha > 0, \quad \Lambda^2 > 0
\end{equation}

ensuring absence of tachyonic instabilities in the scalar sector.

\section{Predictions}

The framework predicts:

\begin{itemize}
\item deviations from GR in strong gravity regimes
\item regularized black hole cores
\item an additional scalar mode
\item modified early-universe dynamics
\end{itemize}

\section{Discussion}

The presented model is an effective truncation of a more general spectral theory. The full structure may include additional curvature invariants and interactions.

\section{Conclusion}

We have presented TIG as a spectral-motivated effective extension of GR. The model introduces curvature-dependent feedback that suppresses divergences and modifies high-energy behavior while recovering GR in the low-energy limit.

\end{document}
